\subsection{An exact counting alternative when few source groups remain}
\label{sec:kirchhoff-sieve}

The common-basis sieve detects a mixed-sign polynomial by randomized
evaluation; distinct copy variables prevent formal cancellation between
different encodings.  When the number $q$ of fixed
mandatory source groups is small, these requirements can instead be
included directly by inclusion--exclusion.  The spanning-tree polynomial
then has positive coefficients, permitting deterministic exact counting.

\begin{theorem}[Kirchhoff counting sieve]\label{thm:kirchhoff-sieve}
Let $q$ be the number of fixed mandatory source SCC groups and $J$ the
number of ideals of the nonfree special-SCC implication poset.  The number
of spanning trees that leave a completely peelable matrix can be computed
exactly and deterministically in
\[
                       2^q J\operatorname{poly}(N)
\]
bit operations.  If this number is positive, a spanning tree and verified
peel order can be extracted within the same asymptotic bound.  In particular,
when $q=0$ the $J\operatorname{poly}(N)$ discovery bound is deterministic.
\end{theorem}
\begin{proof}
Use the original graph edges, with one variable $x_e$ per edge and no
representative/filler copies.  Let $G$ be the reduced oriented incidence
matrix of the graph.  For an ideal $A$ and a subset
$B\subseteq\{1,\ldots,q\}$, ban every edge whose special color mask meets
$A$, and every edge belonging to a mandatory group $S_i$ with $i\in B$.
Let $X_{A,B}(x)$ be the diagonal matrix that retains $x_e$ on the remaining
edges and has zero entries on banned edges.  Put
\[
    R(x)=\sum_{A\in\mathcal I(\mathcal P)}c(A)
             \sum_{B\subseteq[q]}(-1)^{|B|}
                        \det\bigl(GX_{A,B}(x)G^{\mathsf T}\bigr).
\]
Cauchy--Binet gives the spanning-tree form of the matrix--tree theorem
\cite{Chaiken1982}:
\[
        \det\bigl(GXG^{\mathsf T}\bigr)
           =\sum_{T\text{ a spanning tree}}\prod_{e\in T}x_e,
\]
with the sum restricted to retained edges.  A tree minor has determinant
$\pm1$, so the product of its two identical minors is one.  Thus every
original spanning tree occurs with coefficient one, without competing
signed representative encodings.

Inclusion--exclusion over $B$ removes precisely the trees missing a
mandatory group.  The ideal-lattice inversion over $A$ multiplies every
remaining tree coefficient by $f(H(T))$, exactly as in
Theorem~\ref{thm:ideal-sieve}.  Consequently
\[
               R(x)=\sum_{T\text{ successful}}\prod_{e\in T}x_e.
\]
Its value at $x_e=1$ for every edge is the exact nonnegative integer
number of successful trees.

All computations have polynomial bit lengths.  Each unweighted graph
minor counts at most $\binom{m}{r}\leq(m+1)^r$ trees, and
$|c(A)|\leq2^{|\max(A)|}\leq2^d$.  Hence every partial sum in the displayed
formula has bit length at most
\[
             O\bigl(r\log(m+1)+q+\log J+d+1\bigr),
\]
which is polynomial in the explicit input length.  Integer determinants
can be evaluated by fraction-free elimination in polynomial bit time
\cite{Bareiss1968}.
There are at most $2^qJ$ such determinants per counting query.

For witness extraction, scan the graph edges and remove $e$ whenever the
count on the remaining graph stays positive.  If removal fails, every
successful tree in the current graph contains $e$; later removals cannot
invalidate that necessity.  One pass therefore suffices.  At its end the
remaining graph contains a successful tree, and every remaining edge
belongs to every such tree.  Its edge set is thus itself one successful
spanning tree.  There are at most $m$ additional count queries.  Replay
unit-column peeling in the original matrix to produce the final certificate.
\end{proof}

This is an additional exact algorithm for a particular parameter regime,
using classical matrix--tree positivity.  It does not remove the factor
$2^q$ when many independent mandatory source groups remain.  In that
regime, the common-basis construction handles the groups in polynomial
cost, or the deterministic width algorithm may be preferable.
